\documentclass[14pt,a4paper]{extarticle}
\usepackage{fontspec}
\usepackage{polyglossia}
\setmainlanguage{russian}
\setmainfont{Times New Roman}
\usepackage{setspace}
\onehalfspacing
\usepackage{geometry}
\geometry{left=30mm, right=15mm, top=20mm, bottom=20mm}
\usepackage{indentfirst}
\setlength{\parindent}{1.25cm}
\setlength{\parskip}{0pt}
%\usepackage{url}
\usepackage{xurl}
\urlstyle{same}
\usepackage{hyperref}
\usepackage{booktabs}
\usepackage{ragged2e}
\usepackage{wasysym}
\usepackage{graphicx}
\usepackage{float}
\usepackage{enumitem}
\usepackage{array}
\usepackage{siunitx}
\usepackage{microtype}


\begin{document}
	
	\sloppy
	\thispagestyle{empty}
	\centering{
		МИНИСТЕРСТВО НАУКИ И ВЫСШЕГО ОБРАЗОВАНИЯ\\
		РОССИЙСКОЙ ФЕДЕРАЦИИ\\
		Федеральное государственное автономное образовательное учреждение высшего образования <<Санкт-Петербургский политехнический университет Петра Великого>>\\[10pt]
		Институт компьютерных наук и кибербезопасности\\
		Высшая школа технологий искусственного интеллекта\\
		02.03.01 Математика и компьютерные науки\\
		\vspace{\fill}
		{\large
			Отчёт по дисциплине\\
			{\Large \bfseries Цифровая аналитика}\\
			{\Large <<Симулятор гравитационного взаимодействия>>}\\
			\vspace{\fill}
		}


		{
			\flushright
			{\bfseries Выполнил:}\\
			студент группы 5130201/40003\\
			Джабраилов А.В.\\
			\vspace{2pt}
			\hspace{9cm}Подпись: \hrulefill \\
			\vspace{2pt}
			{\bfseries Проверил:}\\
			к. т. н., директор ВШТИИ\\
			Мулюха В.А.\\
			\vspace{2pt}
			\hspace{9cm}Подпись: \hrulefill \\
			\vspace{3pt}
			\hspace{9cm}<<\rule{1.5cm}{0.4pt}>> \hrulefill\ 2026~г.
		}
		
		\centering{
			Санкт-Петербург, 2026
		}
	}
	
	\justifying
	
	\newpage
	\sloppy

    \tableofcontents

    \newpage
  
	\section*{Введение}
	\addcontentsline{toc}{section}{Введение}
	
	Компьютерное моделирование гравитационного взаимодействия используется для того, чтобы наглядно изучать движение тел, устойчивость орбит, столкновения и разрушение объектов под действием приливных сил. Даже простая модель, основанная на законе всемирного тяготения Ньютона, позволяет увидеть сложное поведение системы: сближение тел, образование орбит, хаотические траектории и резкие изменения после столкновений.
	
	Цель проекта --- разработать компьютерное приложение для визуализации эволюции системы космических объектов. Программа моделирует сферические твердые тела в единицах СИ, показывает их движение на трех ортогональных 2D-проекциях и позволяет наблюдать основные физические события: гравитационное движение, столкновения, слияние, фрагментацию и разрушение внутри предела Роша.
	
	Если описать работу симулятора на пользовательском уровне, то система состоит из двух связанных частей. Первая часть выполняет физический расчет: вычисляет ускорения, интегрирует движение и обрабатывает события. Вторая часть отвечает за интерфейс: таблицу объектов, управление временем, загрузку сценариев и визуализацию тел в плоскостях XY, XZ и YZ.
	
	Ключевые задачи проекта:
	\begin{itemize}
		\item реализовать модель ньютоновской гравитации для набора сферических тел;
		\item использовать интегратор Верле для пошагового обновления координат и скоростей;
		\item обработать столкновения твердых тел через слияние или распад на фрагменты;
		\item реализовать разрушение по пределу Роша с условием пребывания внутри зоны не менее 24 часов;
		\item сделать графический интерфейс с тремя ортогональными проекциями и таблицей параметров;
		\item добавить загрузку и сохранение сценариев в CSV-формате;
		\item обеспечить валидацию пользовательских данных и ограничений из технического задания.
	\end{itemize}
	
	Проект реализован на Python. Для графического интерфейса используется PySide6, для двумерной визуализации --- pyqtgraph, для численных расчетов и векторных операций --- NumPy. Разработка велась в формате AI-first: изменения в кодовой базе выполнялись через ИИ-инструменты, а ход работы фиксировался в промптах и файлах проекта.
	
	\newpage
	
	\section{Физическая модель и правила симуляции}
	
	В основе программы лежит модель точечных масс, представленных как однородные сферические тела. Каждое тело имеет имя, массу, радиус, координаты, скорость и ускорение. Все величины задаются в системе СИ.%: масса в килограммах, расстояние в метрах, время в секундах, скорость в метрах в секундах, ускорение в метрах в квадрат секунды.
	
	Гравитационное взаимодействие рассчитывается попарно для всех объектов. Для каждой пары тел определяется вектор между ними, расстояние и соответствующий вклад в ускорение. Такой подход прост и прозрачен: он напрямую повторяет физическую постановку задачи, хотя имеет квадратичную сложность по числу объектов.
	
	Движение тел обновляется интегратором Velocity Verlet. Он удобен для задач небесной механики, потому что использует ускорение на текущем и следующем шаге. На каждом шаге симуляции выполняется следующая последовательность:
	\begin{itemize}
		%\item проверяются настройки фрагментации;
		\item вычисляется шаг времени (берётся из пользовательского интерфейса);
		\item обновляются координаты, скорости и ускорения тел;
		\item обрабатываются столкновения;
		\item проверяется условие предела Роша;
		\item пересчитываются ускорения для отображения в таблице.
	\end{itemize}
	
	Столкновение определяется геометрически: если расстояние между центрами двух тел меньше или равно сумме их радиусов, тела считаются соприкоснувшимися. Далее результат выбирается по правилам технического задания. При очень малом отношении масс происходит слияние. При промежуточном отношении масс используется вероятностный выбор между поглощением и разрушением. Для сопоставимых масс учитывается относительная скорость: если она больше второй космической скорости меньшего тела, происходит фрагментация, иначе тела сливаются.
	
	Слияние реализовано как абсолютно неупругое столкновение. Масса нового тела равна сумме масс, положение и скорость рассчитываются с сохранением импульса, а радиус определяется через сумму объемов. %Это позволяет получить физически понятный результат: два объекта превращаются в один более массивный объект, продолжающий движение с учетом общего импульса системы.
	
	Фрагментация создает набор сферических осколков. Программа сохраняет суммарную массу родителя: масса каждого фрагмента равна массе родительского тела, деленной на фактическое число осколков. Осколки наследуют имя родителя с индивидуальным номером и больше не могут разрушаться повторно. Чтобы облако фрагментов не появлялось в одной точке, центры осколков размещаются случайно внутри объема родительского тела с проверкой на отсутствие пересечений.
	
	Предел Роша используется для моделирования разрушения тела рядом с более массивным объектом. Программа считает плотности первичного тела и спутника, вычисляет расстояние предела Роша и накапливает время пребывания спутника внутри этой зоны. Разрушение происходит только после того, как объект находился внутри предела не менее 24 часов суммарного симулируемого времени.
	
	\newpage
	
	\section{Особенности реализации}
	
	Архитектура проекта разделена на несколько логических слоев. В слое core находятся основные структуры данных: тело, настройки симуляции, системное состояние, константы и правила валидации. Слой physics отвечает за расчет гравитации, интегратор, столкновения, фрагментацию и предел Роша. Слой io реализует загрузку и сохранение CSV-сценариев. Слой ui содержит графический интерфейс, таблицу объектов, панель управления и виджеты проекций.
	
	Центральным объектом физической модели является тело. Оно хранит массу, радиус, положение, скорость, ускорение, признак фрагмента, цвет и визуальную текстуру. %Для безопасной работы с NumPy-массивами используется копирование объектов при передаче между состояниями симуляции. Это важно, потому что координаты, скорости и ускорения изменяются в процессе расчета.
	
	Настройки симуляции включают шаг времени, частоту обновления симуляции, число фрагментов при столкновениях, число фрагментов при разрушении по Рошу и максимальное число объектов. Благодаря этому пользователь может менять скорость моделирования и управлять детализацией разрушения без изменения исходного кода.
	
	Особое внимание уделено валидации данных. Масса и радиус проверяются на попадание в допустимые границы. Координаты, скорости и ускорения должны быть конечными вещественными числами. Имена объектов не должны быть пустыми и не должны повторяться. Также программа запрещает очевидные сложные объекты по имени, например звезды, черные дыры и газовые гиганты, поскольку техническое задание ограничивает модель твердыми телами.
	
	Формат сценариев основан на CSV-файлах. Для каждого объекта в файле задаются имя, масса, радиус, три координаты положения, три компоненты скорости и три компоненты ускорения. Такой формат удобен для ручного редактирования, тестирования и подготовки демонстрационных систем. Встроенные preset-сценарии позволяют быстро запускать примеры вроде системы Земля--Луна, столкновения двух тел, трехтелесной задачи и демонстрации предела Роша.
	
	Графический интерфейс реализован на PySide6. Главное окно состоит из таблицы объектов, панели управления и трех проекций. Таблица позволяет просматривать и редактировать параметры тел, пока симуляция остановлена. Во время запуска редактирование отключается, чтобы пользователь не изменил состояние в середине шага расчета. Панель управления содержит кнопки запуска, паузы и сброса, настройки времени, числа фрагментов, лимита объектов, выбор preset-сценария и дополнительные визуальные переключатели.
	
	Для визуализации используются три ортогональные проекции: XY, XZ и YZ. Тела отображаются окружностями с диаметром, соответствующим физическому радиусу в координатах сцены. Это важно: визуальный контакт тел соответствует тому же условию, которое используется в обработчике столкновений. Дополнительно программа поддерживает цвета, текстуры и подписи объектов, что делает демонстрацию более понятной.
	
	Качество реализации поддерживается тестами. Набор pytest-проверок покрывает валидацию, гравитацию, интегратор Верле, столкновения, фрагментацию, предел Роша, CSV-ввод-вывод, цвета, текстуры и часть поведения интерфейсных моделей. Это помогает безопасно менять физические правила и быстро находить регрессии.
	
	\newpage
	
	\section{Пользовательские сценарии и практическая ценность}
	
	С точки зрения пользователя программа работает как интерактивная лаборатория для изучения гравитационных систем. Можно открыть готовый сценарий, запустить симуляцию и наблюдать, как объекты движутся под действием взаимного притяжения. Три проекции позволяют понять положение тел в пространстве без полноценного 3D-рендеринга.
	
	Основной сценарий использования начинается с выбора preset-сценария или создания собственного набора тел. Пользователь задает массу, радиус, координаты и скорость объекта, после чего запускает симуляцию. Если параметры подобраны для устойчивой орбиты, можно наблюдать движение спутника вокруг более массивного тела. Если траектории пересекаются, программа показывает столкновение и дальнейший результат: слияние или разлет фрагментов.
	
	Отдельный сценарий связан с пределом Роша. Пользователь может загрузить демонстрационный preset, где малое тело находится рядом с массивным объектом, и проследить накопление времени внутри критической зоны. После 24 часов симулируемого времени тело распадается на фрагменты. Такой пример полезен для объяснения того, почему спутники и малые тела могут разрушаться вблизи массивных планет.
	
	Практическая ценность проекта состоит не только в визуализации, но и в возможности экспериментировать с параметрами. Пользователь может менять:
	\begin{itemize}
		\item массы и радиусы тел;
		\item начальные координаты и скорости;
		\item скорость хода симуляции;
		\item число фрагментов при разных типах разрушения;
		\item максимальное число объектов в системе;
		\item готовые сценарии и CSV-файлы.
	\end{itemize}
	
	Для учебных целей программа удобна тем, что связывает формулы с наблюдаемым поведением. Пользователь видит не только таблицу чисел, но и траектории, сближения, столкновения и последствия разрушения. Это снижает порог входа в тему численного моделирования и помогает быстрее понять, как параметры влияют на движение системы.
	
	Для команды проекта важна и инженерная ценность. Кодовая база разделена на независимые модули, поэтому физику, интерфейс и ввод-вывод можно развивать отдельно. Наличие тестов и CSV-сценариев упрощает проверку новых идей. Такой проект можно использовать как основу для дальнейших экспериментов: повышения производительности, добавления новых сценариев, улучшения визуализации или сравнения разных численных методов.
	
	\newpage
	
	\section*{Заключение}
	\addcontentsline{toc}{section}{Заключение}
	
	В рамках проекта был реализован симулятор гравитационного взаимодействия сферических твердых тел. Программа объединяет физический расчет, обработку событий и графический интерфейс, позволяя наблюдать эволюцию системы в трех ортогональных проекциях и одновременно контролировать численные параметры через таблицу.
	
	Ключевой результат работы --- цельное приложение, в котором реализованы:
	\begin{itemize}
		\item ньютоновская гравитация между телами;
		\item интегратор Velocity Verlet;
		\item обработка столкновений через слияние и фрагментацию;
		\item разрушение по пределу Роша после 24 часов пребывания внутри зоны;
		\item CSV-загрузка и сохранение сценариев;
		\item графический интерфейс с управлением симуляцией и тремя 2D-проекциями;
		\item валидация пользовательских данных и набор автоматических тестов.
	\end{itemize}
	
	Проект показывает, что даже в рамках ограниченной учебной модели можно получить наглядный цифровой продукт: пользователь задает начальные условия, запускает расчет и сразу видит последствия физических правил. При этом архитектура остается достаточно простой для сопровождения: основные правила физики отделены от интерфейса, а сценарии можно хранить и переносить в CSV-файлах.
	
	Отдельная особенность работы --- формат AI-first. Код и документация развивались итерационно через постановку задач ИИ-ассистенту, проверку результата и уточнение требований. Такой процесс позволил быстро проходить путь от базовой симуляции к более сложным сценариям: фрагментации, пределу Роша, визуальным улучшениям и дополнительным проверкам корректности.
	
	В дальнейшем проект можно развивать в нескольких направлениях: оптимизировать расчет для большого числа тел, расширить набор демонстрационных сценариев, добавить более подробную диагностику физических событий и улучшить визуальную часть интерфейса. Уже текущая версия решает основную задачу: делает гравитационную динамику, столкновения и разрушение тел понятными и наблюдаемыми в интерактивном режиме.
	
\end{document}
